\documentclass[12pt]{article}
% \usepackage[a4paper, margin=1in]{geometry}
% \usepackage{graphicx}
\input preamble.tex
\input macros.tex

\usepackage[most]{tcolorbox}
\usepackage{enumitem}
\usepackage{wrapfig}
\usepackage{bm}
\usepackage{stmaryrd}
\usepackage{caption}
\usepackage{subcaption}

\usepackage{longtable}

\title{Understanding the loss landscape of deep neural networks: new ideas}
%\author{[Organization Name / Logo]}
\date{}

\begin{document}

\maketitle

\begin{abstract}
We study the loss landscape of deep neural networks with an emphasis on higher-order geometry and the behavior of stochastic gradient descent (SGD) in flat or degenerate regions. We take two complementary viewpoints: a global one via convergence guarantees under regularity assumptions, and a local one via spectral statistics of curvature. Our main observation is that lazy training can be interpreted as a form of convexification of the landscape, which clarifies when higher-order information matters. We conclude with research directions and leave empirical plots to Appendix~\ref{app:plots}.
\end{abstract}

\tableofcontents

\newpage

\section{Introduction}
In the past few years, research has focused on understanding why over\-parameterized deep networks train efficiently despite highly non-convex objectives. Central themes include the structure of critical points, implicit bias of optimization, and properties of curvature along training trajectories. It remains unclear why simple first-order methods work so well in regimes where the Hessian is often degenerate and flat directions abound.

This work aims to clarify how higher-order geometry interacts with SGD noise and step size, and to delineate regimes where local convexification occurs (e.g., lazy training). Our principal findings are outlined in the abstract; the paper is structured as follows: Section~\ref{sec:lit} reviews the literature, Section~\ref{sec:directions} proposes concrete research directions, and Appendix~\ref{app:plots} gathers plots.

\section{Literature review}\label{sec:lit}
We are compiling a BibTeX corpus covering optimization theory for deep networks, spectral analyses of Hessians, and connections to generalization. To understand how SGD behaves in flat or degenerate landscapes, higher-order derivatives become crucial. A simple curvature summary that often appears in analyses is
\begin{equation}\label{eq:scalar-curvature}
\bm{S} \;=\; \operatorname{tr}\!\big(\bm{H}\big)^{2} \;-\; \operatorname{tr}\!\big(\bm{H}^{2}\big),
\end{equation}
where \(\bm{H}\) denotes the Hessian and \(\operatorname{tr}(\cdot)\) is the trace operator. Equation~\eqref{eq:scalar-curvature} highlights how the spectrum of \(\bm{H}\) controls aggregated curvature.

Broadly, there are two paths to understand optimization:
\paragraph{Global analysis.} One seeks convergence guarantees under regularity assumptions (e.g., sufficiently small learning rate, smoothness, and controlled non-convexity measured via Hessian bounds). This perspective prioritizes end-to-end guarantees.
\paragraph{Local spectral analysis.} One studies local curvature statistics (eigenvalue distributions, outlier spikes, anisotropy) and their interaction with step size and SGD noise. This perspective prioritizes mechanistic understanding near iterates.

A useful heuristic is that lazy training acts as a \emph{convexification} of the landscape around initialization, which can be diagnosed via spectra of \(\bm{H}\) and the stability of the neural tangent kernel. We defer empirical illustrations to Appendix~\ref{app:plots}.

\section{Which direction to work on}\label{sec:directions}
We focus on formalizing when local convexification emerges and how it depends on architecture, width, depth, and initialization scale. First, we will characterize regimes where higher-order terms beyond the quadratic approximation materially affect descent dynamics. Second, we will connect spectral signals in \(\bm{H}\) to the effective noise of SGD and to step-size schedules. Finally, we will test whether lazy training reliably predicts optimization stability across models and datasets.

\section{Conclusion}
We outlined two complementary approaches to the loss landscape---global guarantees and local spectral analysis---and argued that lazy training can be read as convexification. Limitations include the need for stronger empirical coverage and tighter non-asymptotic bounds. We recommend a combined program: sharpen spectral diagnostics, couple them with SGD noise models, and validate across scales. All plots are collected in Appendix~\ref{app:plots}.

\appendix
\section{Additional plots}\label{app:plots}
All plots referenced in the main text are reported here with captions and cross-references. We keep the main exposition uncluttered and refer to the appendix for visual summaries and ablations.

\end{document}